\section{Redistricting as MCMC on a Graph}
\label{sec:mcmc}

\subsection{State Space}

Let $G = (V, E)$ be the census-tract adjacency graph for a state, with
$|V| = n$ tracts.
A \emph{redistricting plan} is a surjection $\pi : V \to [k]$ such that:
\begin{enumerate}
\item \textbf{Population balance}: for each district $i$,
      $|P(\pi^{-1}(i)) - P_{\rm target}| \leq \epsilon \cdot P_{\rm target}$,
      where $P(S)$ is the total population of tracts in $S$ and
      $\epsilon = 0.005$ (the DIA tolerance).
\item \textbf{Contiguity}: for each district $i$, the induced subgraph
      $G[\pi^{-1}(i)]$ is connected.
\end{enumerate}

Let $\mathcal{P}$ denote the set of valid plans.
$|\mathcal{P}|$ is exponentially large in $n$ (lower bounded by
$(n/k)^k$ for balanced partitions~\citep{duchindist2020}).

\subsection{The ReCom Proposal Distribution}

The \recom{} chain~\citep{deford2021} transitions from plan $\pi$ to
plan $\pi'$ via:
\begin{enumerate}
\item Select an edge $(u, v) \in E$ such that $\pi(u) \neq \pi(v)$
      (a \emph{boundary edge}).
\item Let $A = \pi^{-1}(\pi(u))$ and $B = \pi^{-1}(\pi(v))$.
\item Construct the merged subgraph $G[A \cup B]$.
\item Sample a random spanning tree $\mathcal{T}$ of $G[A \cup B]$.
\item Remove from $\mathcal{T}$ the unique edge $(s, t)$ such that the
      resulting bipartition satisfies population balance.
\item Assign $\pi'(v) = \pi(u)$ for all $v$ in one component and
      $\pi'(v) = \pi(v)$ for the other.
\end{enumerate}

This defines a Markov chain on $\mathcal{P}$.
For the chain to have a well-defined stationary distribution, steps that
produce invalid plans (no balancing edge exists) are rejected; the
chain stays at $\pi$.

\subsection{Stationary Distribution}

Under uniform spanning tree sampling, the \recom{} chain is designed
to have stationary distribution approximately uniform over $\mathcal{P}$
\citep{deford2021}.
The approximation arises because the probability of selecting each
boundary edge is not exactly uniform; a Metropolis correction can make
the chain exactly reversible but at the cost of higher rejection rates.

For the mixing time analysis in Section~\ref{sec:bounds}, we analyze the
idealized lazy reversible version of \recom{} that holds with probability
$1/2$ at each step and otherwise applies the \recom{} proposal with
Metropolis correction.
This chain has a well-defined spectral gap.

\subsection{The Hamming Distance and Plan Space Geometry}

The Hamming distance $d_{\rm Ham}(\pi, \pi') = \frac{1}{n}|\{v : \pi(v) \neq \pi'(v)\}|$
defines a metric on $\mathcal{P}$.
Under \recom{}, each accepted step changes $\approx n/k$ tracts, giving
$d_{\rm Ham}(\pi_t, \pi_{t+1}) \approx 1/k$.
The diameter of $\mathcal{P}$ under Hamming distance is at most 1
($(k-1)/k$ in practice).
For the chain to mix, it must travel a Hamming distance of approximately
$1/2$ from any starting plan.
At $1/k$ per step, this requires $\approx k/2$ steps in the best case---but
the actual mixing time is much larger because the chain must navigate the
complex graph structure of $\mathcal{P}$.
